% Part: lambda-calculus
% Chapter: syntax
% Section: substitution

\documentclass[../../../include/open-logic-section]{subfiles}

\begin{document}

\olfileid{lam}{syn}{sub}
\olsection{Substitusi}

\begin{explain}
Variabel bebas iku rujukan marang variabel
ing lingkungan, mula ana gunane yen kita
nggunakake biji tartamtu kanggo ngganti
variabel bebas. Umpamane, kita bisa
ngganti $f$ ing $\lambd[x][f x]$
nganggo suku tartamtu, kayata fungsi
identitas $\lambd[y][y]$. Asile
$\lambd[x][(\lambd[y][y]) x]$.
Proses ngganti variabel bebas nganggo
suku lambda diarani substitusi.
\end{explain}

\begin{defn}[Substitusi] \ollabel{defn:substitution}
  \emph{Substitusi} suku $N$ kanggo
  variabel $x$ ing suku $M$,
  $\Subst{M}{N}{x}$, ditetepake
  kanthi induktif kaya mangkene:
  \begin{enumerate}
    \item $\Subst{x}{N}{x} = N$. \ollabel{defn:substitution-1} 
    \item $\Subst{y}{N}{x}  = y$ yen $x \neq y$. \ollabel{defn:substitution-2} 
    \item $\Subst{PQ}{N}{x}  = (\Subst{P}{N}{x}) (\Subst{Q}{N}{x})$.
      \ollabel{def:substitution-3}
    \item $\Subst{(\lambd[y][P])}{N}{x} = \lambd[y][\Subst{P}{N}{x}]$,
      yen $x \neq y$ lan $y \notin \FV{N}$.
      \ollabel{defn:substitution-4}
    \item $\Subst{(\lambd[x][P])}{N}{x} = \lambd[x][P]$.
    \item Yen $x \notin \FV{P}$,
      $\Subst{(\lambd[y][P])}{N}{x} = \lambd[y][P]$.
  \end{enumerate}
  Ing kasus abstraksi liyane kang ora nyukupi
  syarat kasebut, substitusi ora katetepake.
\end{defn}

\begin{explain}
Ing \olref{defn:substitution}\olref{defn:substitution-4},
syarat $x \neq y$ njaga supaya
muncule variabel~$x$ kang \emph{kaiket}
ing~$M$ ora diganti dening~$N$.
Umpamane, substitusi
$\Subst{\lambd[x][x]}{y}{x}$
asile dudu $\lambd[x][y]$,
nanging tetep $\lambd[x][x]$.

Nalika nyubstitusi $N$ kanggo~$x$ ing
$\lambd[y][P]$, kita uga mbutuhake
$y \notin \FV{N}$. Umpamane,
kita ora bisa nyubstitusi $y$ kanggo $x$
ing $\lambd[y][x]$, yaiku
$\Subst{\lambd[y][x]}{y}{x}$,
amarga iku bakal ngasilake
$\lambd[y][y]$, suku kang makili
fungsi kang nampa argumen lan
langsung mbalekake argumen iku.
Nanging $\lambd[y][x]$ makili
fungsi kang tansah mbalekake suku~$x$
(utawa apa wae kang dirujuk dening $x$).
Mula asil kang kita karepake yaiku
fungsi kang nampa argumen,
nglirwakake argumen iku, lan
mbalekake variabel lingkungan~$y$.
Kanggo nindakake iki kanthi bener,
kita kudu luwih dhisik ngganti jeneng
variabel~$y$ kang kaiket.
\end{explain}

\begin{prob}
  Apa asil substitusi ing ngisor iki?
  \begin{enumerate}
  \item $\Subst{\lambd[y][x(\lambd[w][vwx])]}{(uv)}{x}$
  \item $\Subst{\lambd[y][x(\lambd[x][x])]}{(\lambd[y][xy])}{x}$
  \item $\Subst{y(\lambd[v][xv])}{(\lambd[y][vy])}{x}$
  \end{enumerate}
\end{prob}

\begin{thm} \ollabel{thm:notinfv}
  Yen $x \notin \FV{M}$, mula
  $\FV{\Subst{M}{N}{x}} = \FV{M}$
  yen substitusine katetepake.
\end{thm}

\begin{proof}
  Kanthi induksi tumrap pambentukan $M$.
  \begin{enumerate}
  \item $M$ iku variabel: latihan.
  \item $M$ awujud $(PQ)$: latihan.
  \item $M$ awujud $\lambd[y][P]$.
    Yen $x=y$ utawa
    $x \notin \FV{P}$, substitusi
    ora ngowahi abstraksi iki,
    mula pratelan cetha.
    Yen $x \neq y$ lan
    aturan rekursif lumaku,
    $\Subst{\lambd[y][P]}{N}{x}$
    katetepake, asile
    $\lambd[y][\Subst{P}{N}{x}]$.
    Mula $\Subst{P}{N}{x}$
    katetepake. Saka
    $x \notin \FV{M}$ lan $x \neq y$
    uga oleh $x \notin \FV{P}$.
    Dadi:
    \begin{multline*}
      \FV{\Subst{\lambd[y][P]}{N}{x}} \\
    \begin{aligned}
      & = \FV{\lambd[y][\Subst{P}{N}{x}]} &&
      \text{miturut \olref{defn:substitution-4}}\\
      & = \FV{\Subst{P}{N}{x}} \setminus \{y\} &&
      \text{miturut \olref[fv]{def:fv}\olref[fv]{def:fv2}}\\
      & = \FV{P} \setminus \{y\} && \text{miturut hipotesis induksi} \\
      & = \FV{\lambd[y][P]} && \text{miturut \olref[fv]{def:fv}\olref[fv]{def:fv2}}
    \end{aligned}
    \end{multline*}
  \end{enumerate}
\end{proof}

\begin{prob}
  Rampungna bukti \olref[lam][syn][sub]{thm:notinfv}.
\end{prob}

\begin{thm} \ollabel{thm:infv}
  Yen $x \in \FV{M}$, mula
  $\FV{\Subst{M}{N}{x}} = (\FV{M} \setminus
  \{x\}) \cup \FV{N}$, yen substitusine
  katetepake.
\end{thm}

\begin{proof}
  Kanthi induksi tumrap pambentukan $M$.
  \begin{enumerate}
  \item $M$ iku variabel: latihan.
  \item $M$ awujud $PQ$: Amarga
    $\Subst{(PQ)}{N}{x}$ katetepake,
    asile mesthi
    $(\Subst{P}{N}{x})(\Subst{Q}{N}{x})$,
    lan substitusi ing loro-lorone
    subsuku katetepake. Saka
    $x \in \FV{PQ}$, oleh
    $x \in \FV{P}$ utawa
    $x \in \FV{Q}$, utawa loro-lorone.
    Sisa buktine dadi latihan.
  \item $M$ awujud $\lambd[y][P]$.
    Amarga $x \in \FV{M}$,
    mesthi $x \neq y$ lan
    $x \in \FV{P}$.
    Amarga
    $\Subst{\lambd[y][P]}{N}{x}$
    katetepake, asile
    $\lambd[y][\Subst{P}{N}{x}]$;
    $\Subst{P}{N}{x}$ katetepake
    lan $y \notin \FV{N}$.
    Mula:
    \begin{multline*}
      \FV{\Subst{(\lambd[y][P])}{N}{x}} \\
      \begin{aligned}
      & = \FV{\lambd[y][\Subst{P}{N}{x}]} \\
      & = \FV{\Subst{P}{N}{x}} \setminus \{y\} \\
      & = ((\FV{P} \setminus \{x\}) \cup \FV{N}) \setminus \{y\}
       && \text{miturut hipotesis induksi} \\
      & = (\FV{P} \setminus \{x, y\}) \cup \FV{N}
       && y \notin \FV{N} \\
      & = (\FV{\lambd[y][P]} \setminus \{x\}) \cup \FV{N}
      \end{aligned}
      \end{multline*}
  \end{enumerate}
\end{proof}

\begin{prob}
  Rampungna bukti \olref[lam][syn][sub]{thm:infv}.
\end{prob}

\begin{thm}\ollabel{thm:clr}
  $x \notin \FV{\Subst{M}{N}{x}}$
  yen substitusine katetepake
  lan $x \notin \FV{N}$.
\end{thm}

\begin{proof}
  Latihan.
\end{proof}

\begin{prob}
  Buktikna \olref[lam][syn][sub]{thm:clr}.
\end{prob}


\begin{thm}\ollabel{thm:inv}
  Yen $\Subst{M}{y}{x}$ katetepake
  lan $y \notin \FV{M}$,
  mula $\Subst{\Subst{M}{y}{x}}{x}{y}$
  katetepake lan
  $\Subst{\Subst{M}{y}{x}}{x}{y} = M$.
\end{thm}

\begin{proof}
  Yen $x=y$, loro substitusi iku
  identitas, mula pratelan cetha.
  Yen $x \notin \FV{M}$, aturan
  substitusi nuduhake yen loro-lorone
  uga ora ngowahi $M$.
  Saiki anggep $x \neq y$ lan
  $x \in \FV{M}$, banjur gunakna
  induksi tumrap pambentukan $M$.
  \begin{enumerate}
  \item $M$ iku variabel $z$: latihan.
  \item $M$ awujud $(PQ)$. Mula:
    \begin{align*}
      \Subst{\Subst{(PQ)}{y}{x}}{x}{y}
      &=\Subst{((\Subst{P}{y}{x})(\Subst{Q}{y}{x}))}{x}{y} \\
      &= (\Subst{\Subst{P}{y}{x}}{x}{y})(\Subst{\Subst{Q}{y}{x}}{x}{y}) \\
      &= (PQ) \text{ miturut hipotesis induksi}
    \end{align*}
  \item $M$ awujud $\lambd[z][N]$.
    Amarga $\Subst{\lambd[z][N]}{y}{x}$
    katetepake lan $x \in \FV{M}$,
    $z \neq x$ lan $z \neq y$.
    Saka $y \notin \FV{M}$ lan
    $z \neq y$, uga oleh
    $y \notin \FV{N}$.
    Miturut hipotesis induksi, substitusi
    balike ing~$N$ katetepake. Mula:
    \begin{align*}
      \Subst{\Subst{(\lambd[z][N])}{y}{x}}{x}{y}\\
      & = \Subst{(\lambd[z][\Subst{N}{y}{x}])}{x}{y} \\
      & = \lambd[z][\Subst{\Subst{N}{y}{x}}{x}{y}] \\
      &= \lambd[z][N] \text{ miturut hipotesis induksi} 
    \end{align*}
  \end{enumerate}
\end{proof}

\begin{prob}
  Rampungna bukti \olref[lam][syn][sub]{thm:inv}.
\end{prob}

\end{document}
